\section{Open Questions：何时检索、怎样扩展、证据是否被使用}

课堂在 Atlas 后进入开放问题，并穿插学生 Q\&A。这里的重点从“哪个架构更强”转向动态 policy、scaling、legal-risk motivation、context position、tool use 与 instruction tuning。很多问题直到今天仍不能用单一 benchmark 回答。

\subsection{RAG 与 Long Context 不是完全相反}

学生问超长 context 是否会取代 RAG。讲者回答，把整本 Harry Potter 放进 context 只为寻找猫头鹰名字非常低效；而长上下文方法为了扩展 attention 往往引入 sparsity，本质上也在选择少量相关连接。

\begin{knowledgebox}{Long context 与 retrieval 的共同问题}
两者都要决定“哪些 token 值得计算”。Long context 把选择放在 attention pattern 或 cache 管理中，RAG 把选择放在 external index；最终都面对 relevance、latency、position bias 与 evidence verification。
\end{knowledgebox}

\teachervoice{课堂判断是：如果长上下文真正扩到巨大规模，它会越来越像 sparse/non-parametric retrieval，而不是无代价地对所有 token 做 full attention。}

\subsection{When to Retrieve：把检索变成策略}

每个 token 都检索会浪费 compute，只在开头检索一次又可能错过后续知识需求。FLARE 让模型根据生成置信度或未来内容主动决定何时搜索、搜索什么，使 retrieval 从固定 pipeline 变成 sequential policy。

\lecturefigure{slide-39-when-to-retrieve.jpg}{FLARE 与 active retrieval：决定何时、为何检索}{Stanford 课堂视频 00:56:36}

\begin{importantbox}{Retrieval policy 的成本目标}
令动作 $a_t\in\{\text{retrieve},\text{continue}\}$，检索成本为 $c_R$，错误成本为 $c_E$，可优化
\[
\mathbb{E}\!\left[\sum_t c_R\mathbf{1}(a_t=\text{retrieve})
+c_E\mathbf{1}(\hat y\ne y)\right].
\]
Policy 要学习在信息价值高于检索成本时调用工具，而不是把更多 retrieval calls 当作天然更好。
\end{importantbox}

\subsection{How to Train at Scale}

Scale slide 讨论 full async update、query-only update、TRIME/近似方法，以及用 BM25 构造相似 batch 后做 in-batch training。Index 越大，document-side full refresh 越贵，因此 hard-negative sampling、memory bank 与 reranker 成为重要近似。

\lecturefigure{slide-40-training-at-scale.jpg}{大规模 retriever training 的 in-batch 与 memory 方法}{Stanford 课堂视频 00:57:39}

\begin{knowledgebox}{读图：In-batch negatives 节省什么}
一个 batch 中其他 documents 可作为 negatives，避免每个 query 单独搜索大量负例；memory bank 扩大负例集合。风险是 false negatives：语义相关文档被当成负例，会扭曲 representation。
\end{knowledgebox}

\subsection{SILO：把 legal-risk motivation 放进架构}

SILO 提出在较安全、许可更清晰的数据上训练 parametric model，再在测试时从 non-parametric datastore 访问更广内容。课堂把它视为应对 copyright provenance 与 legal risk 的优雅方向，但并未声称 index 中的内容因此免于法律责任。

\lecturefigure{slide-41-silo-legal-risk.jpg}{SILO：用 non-parametric datastore 隔离部分训练数据风险}{Stanford 课堂视频 00:58:45}

\begin{warningbox}{Architecture 不会自动解决法律问题}
Datastore 仍需考虑授权、访问、日志、删除、输出引用和用户数据；retrieved content 也可能被模型改写。SILO 的研究价值是把训练数据与 test-time knowledge 分离，便于替换和审计，不是法律结论。
\end{warningbox}

\teachervoice{讲者把 SILO 与当时针对模型训练数据来源的诉讼联系起来，强调 parametric model 可以在 public-domain 数据上训练，再从更高风险 index 检索。讲义保留其 motivation，并明确不把它升级为合规保证。}

\subsection{Lost in the Middle：检索正确也可能使用失败}

Lost-in-the-Middle slide 的曲线显示 relevant document 位于 context 开头或结尾时更易被使用，放在中间性能下降。若 generator 完全遵守 retriever ranking，位置变化不应如此显著；这证明 evidence placement 与 attention behavior 是独立 failure mode。

\lecturefigure{slide-42-lost-in-middle.jpg}{Relevant passage 在长 context 中间时容易被忽略}{Stanford 课堂视频 01:00:45}

\begin{knowledgebox}{读图：U-shaped curve 说明什么}
横轴是 relevant document 在 20 个 retrieved documents 中的位置，纵轴是任务表现。开头/结尾高、中间低，说明模型具有 position bias。它不证明中间文档永远无效，而是要求 reranking、context packing 与 evaluation 同时考虑位置。
\end{knowledgebox}

\begin{warningbox}{Retrieval recall 与 answer grounding 要分层测}
可记录：关键 evidence 是否进入 top-$k$、是否进入最终 prompt、位于什么位置、答案是否引用它、输出是否与 evidence 一致。只看 final exact match 无法定位是 search、packing 还是 generation 失败。
\end{warningbox}

\subsection{Toolformer：Retrieval 是更一般的 Tool Use}

Toolformer slide 把 search API 放在 calculator、calendar、translation 等工具集合中。Language model 可以学习何时插入 API call、读取返回值并继续生成；retrieval 因而是 tool augmentation 的一个实例。

\lecturefigure{slide-43-toolformer.jpg}{Toolformer：语言模型自监督学习调用外部工具}{Stanford 课堂视频 01:01:57}

\begin{knowledgebox}{读图：从 RAG 到 Agent 的接口变化}
固定 RAG 每个请求都调用检索；tool-using LM 可以选择 search、calculator 或其他 API。系统新增 action selection、argument generation、tool error handling 与 permission boundary，问题从静态 architecture 扩展为 sequential decision making。
\end{knowledgebox}

\teachervoice{讲者说，active web search 不一定访问自有 index；更一般地，LM 是 tool user，retrieval 只是众多工具之一，而真正学习动作选择可能需要 RL。}

\subsection{Self-RAG：Retrieve、Generate、Critique}

Self-RAG 让单一 LM 生成 reflection tokens，决定是否检索、评估 evidence relevance、生成答案并 critique support/utility。它试图把 active retrieval 与自我评价纳入同一训练框架。本节承接 Toolformer 的动作选择，并进一步追问同一个模型如何判断证据质量与答案支持度。

\lecturefigure{slide-44-self-rag.jpg}{Self-RAG 的主动检索与 reflection 流程}{Stanford 课堂视频 01:02:15}

\begin{importantbox}{Self-reflection 仍需外部评测}
Critique token 是模型预测，不是事实 oracle。系统应分别验证 retrieval decision、source relevance、claim support 与 final utility；模型自评可作为 feature，却不能取代 held-out labels、human review 或可执行 verifier。
\end{importantbox}

\begin{lstlisting}[caption={Active retrieval 与 evidence verification 循环}]
state = initialize(question)
while not state.finished and state.steps < max_steps:
    if policy.needs_evidence(state):
        docs = search(state.query)
        docs = rerank_and_filter(docs, source_policy)
        state.add_evidence(docs)
    draft = generator.continue_answer(state)
    verdict = verify_claims(draft, state.evidence)
    state = revise_or_finish(state, draft, verdict)
return state.answer
\end{lstlisting}

\subsection{Instruction Tuning the Entire RAG System}

Instruction-tuning slide 对比 InstructRetro 与 RA-DIT：过去指令数据主要训练 generator，retriever 不一定理解“比较、引用、只使用指定来源”等任务意图。Dual instruction tuning 同时适配 retrieval 与 generation，使整套系统遵循 instruction。

\lecturefigure{slide-45-instruction-tuning.jpg}{InstructRetro 与 RA-DIT：对 RAG system 做 instruction tuning}{Stanford 课堂视频 01:03:03}

\begin{knowledgebox}{Instruction 对 Retriever 也有语义}
``Give a counterargument''、``use peer-reviewed sources'' 与 ``summarize this document'' 需要不同 evidence。若 retriever 只看 topic similarity，可能返回主题相关却不满足 task operator 的文档；instruction-aware query representation 可减少这种错配。
\end{knowledgebox}

\subsection{Advanced Frozen RAG：Developer Community 的实用创新}

即使不做 joint training，开发者也构建了 child-parent recursive retriever、hybrid search/RRF、zero-shot LLM reranker 与 HyDE。HyDE 先让模型生成 hypothetical document，再用它的 embedding 搜真实文档，课堂称其为“用 hallucination 修复 hallucination”的有趣想法。

\lecturefigure{slide-46-advanced-frozen-rag.jpg}{Advanced Frozen RAG 的层级检索、融合、reranking 与 HyDE}{Stanford 课堂视频 01:04:30}

\begin{knowledgebox}{术语消化：四个 advanced frozen patterns}
\begin{tabularx}{\textwidth}{L{0.25\textwidth}>{\raggedright\arraybackslash}X}
\toprule
模式 & 机制 \\
\midrule
Child-parent retrieval & 用小 chunk 精确召回，再扩展到父段落提供完整 context \\
Hybrid/RRF & 合并 BM25 与 dense ranking，兼顾 exact 和 semantic match \\
LLM reranker & 让强 LM 在小候选集上比较 relevance，成本较高 \\
HyDE & 先生成 hypothetical answer/document，再用其 embedding 检索真实证据 \\
\bottomrule
\end{tabularx}
\end{knowledgebox}

\begin{warningbox}{HyDE 的 hypothetical text 不是证据}
它只用来构造更丰富的 query representation；最终回答仍应引用真实 retrieved documents。若直接把 hypothetical document 当事实，会把 retrieval query expansion 变成自证循环。
\end{warningbox}

\teachervoice{讲者一方面批评 frozen RAG 的 Frankenstein 属性，另一方面肯定 LlamaIndex、LangChain、Chroma、Weaviate 等开发者生态做出的实用创新。研究方向与工程 baseline 可以同时成立。}

\subsection{本章小结}

Open questions 集中在 policy 与系统边界：何时检索、怎样扩展 retriever、如何处理 legal-risk motivation、怎样保证证据真正被 generator 使用，以及 retrieval 如何成为 tool use。Frozen RAG 仍有强工程技巧，但 joint learning 与可分层 evaluation 决定其上限。

